\documentclass[12pt,a4paper]{article}
\usepackage[margin=25mm, headheight=15pt, headsep=10pt]{geometry}
\usepackage{fontspec}
\usepackage[table]{xcolor} % Note: table option is required for rowcolors
\usepackage{titlesec}
\usepackage{tocloft}
\usepackage{fancyhdr}
\usepackage{booktabs}
\usepackage{tcolorbox}
\tcbuselibrary{skins, breakable}
\usepackage[colorlinks=true, linkcolor=emerald, urlcolor=emerald, bookmarks=true]{hyperref}

\definecolor{emerald}{RGB}{0,102,51}
\definecolor{darkred}{RGB}{139,0,0}
\definecolor{lightgray}{RGB}{245,245,245}

\usepackage[nil]{babel}
\babelprovide[import, main]{bengali}
\babelprovide[import]{arabic}
\babelfont{rm}[Renderer=HarfBuzz,Scale=1.0]{Noto Serif Bengali}
\babelfont{sf}[Renderer=HarfBuzz]{Noto Sans Bengali}
\babelfont{tt}[Renderer=HarfBuzz]{Noto Sans Bengali}
\babelfont[arabic]{rm}[Renderer=HarfBuzz,Scale=1.1]{Noto Naskh Arabic}

% Section styling
\titleformat{\section}{\Large\bfseries\color{emerald}}{\thesection.}{0.5em}{}
\titleformat{\subsection}{\large\bfseries\color{emerald!85!black}}{\thesubsection.}{0.5em}{}
\titleformat{\subsubsection}{\normalsize\bfseries\color{emerald!70!black}}{\thesubsubsection.}{0.5em}{}
\titlespacing*{\section}{0pt}{18pt}{8pt}

% Fancy Header/Footer setup
\pagestyle{fancy}
\fancyhf{} % clear all header and footer fields
\fancyhead[L]{\color{emerald}\small\textbf{\nouppercase{\leftmark}}}
\fancyfoot[L]{\scriptsize\color{black!50}{\lat{AI}}-সংকলিত, আলেম-যাচাইকৃত নয়}
\fancyfoot[C]{\color{emerald!70!black}\thepage}
\renewcommand{\headrulewidth}{0.4pt}
\renewcommand{\headrule}{\hbox to\headwidth{\color{emerald}\leaders\hrule height \headrulewidth\hfill}}
\renewcommand{\footrulewidth}{0pt}

% Custom boxes
% 1. Ayat/Hadith Box
\newtcolorbox{ayatbox}[1][]{
    enhanced, breakable, 
    colback=emerald!5, 
    colframe=emerald, 
    boxrule=0.5pt, 
    arc=3pt, 
    left=10pt, right=10pt, top=10pt, bottom=10pt,
    #1
}

% 2. Quote/Translation Box
\newtcolorbox{quotebox}[1][]{
    enhanced, breakable,
    frame hidden,
    colback=lightgray,
    borderline west={3pt}{0pt}{emerald},
    left=15pt, right=10pt, top=10pt, bottom=10pt,
    sharp corners,
    #1
}

% 3. AI-generated notice box
\newtcolorbox{warnbox}[1][]{
    enhanced, breakable,
    colback=red!4,
    colframe=darkred,
    boxrule=0.8pt,
    arc=3pt,
    left=10pt, right=10pt, top=10pt, bottom=10pt,
    #1
}

% Commands
\newcommand{\arab}[1]{\foreignlanguage{arabic}{#1}}
\newcommand{\kib}[1]{\textcolor{emerald}{\textbf{#1}}}

% Latin (English) fallback: Noto Serif Bengali has NO Latin glyphs
\newfontfamily\latfont[Renderer=HarfBuzz]{Noto Serif}
\newcommand{\lat}[1]{{\latfont #1}}

\setlength{\parskip}{5pt}
\setlength{\parindent}{0pt}

% Fix for missing bullet in Noto Serif Bengali
\renewcommand{\labelitemi}{$\bullet$}



\begin{document}
\begin{center}{\Large\bfseries\color{emerald} ক্ষমা করা পাপ — কিয়ামতের দিন কী হয়? (কুরআনের আয়াত)}\end{center}
\vspace{1em}
\section*{ক্ষমা করা পাপ — কিয়ামতের দিন কী হয়? (কুরআনের আয়াত)}

\begin{warnbox}
 \textbf{গুরুত্বপূর্ণ নোটিশ (\lat{Important} \lat{Notice}):} এই কন্টেন্টটি \lat{AI} (কৃত্রিম বুদ্ধিমত্তা) দিয়ে প্রস্তুত ও সংকলিত। \lat{AI} কঠোর সূত্র-মান নীতি মেনে শুধু নির্ভরযোগ্য উৎস থেকে তথ্য সংগ্রহ করেছে — কেবল সহীহ/হাসান হাদিস ও সহীহ সনদের আসার রাখা হয়েছে, দুর্বল সনদ বাদ দেওয়া হয়েছে। তবে এটি কোনো আলেম বা মুহাদ্দিস কর্তৃক যাচাইকৃত নয়।
\end{warnbox}

\begin{quote}
\textbf{পড়ার নিয়ম:} আয়াতগুলো ৪টি থিমে সাজানো। প্রতিটিতে — সূত্র কার্ড $\rightarrow$ মূল আরবি (যাচাইকৃত) $\rightarrow$ উচ্চারণ $\rightarrow$ বাংলা অর্থ (\lat{pure}) $\rightarrow$ শব্দের ব্যাখ্যা $\rightarrow$ সংক্ষিপ্ত তাফসীর (ইবনে কাসীর, সা'দী প্রভৃতি থেকে; পৃষ্ঠা নম্বর ইচ্ছাকৃতভাবে দেওয়া হয়নি — উৎস গ্রন্থ নাম উল্লেখ করা হলো) $\rightarrow$ সংশ্লিষ্ট দলিল। \\
\textbf{মূলনীতি:} শুধু সহীহ/হাসান হাদিস; দুর্বল সনদ বাদ।
\end{quote}

\medskip\hrule\medskip

\subsection*{থিম ১ — তাওবার পর পাপের অবস্থা: মুছে যাওয়া ও নেকিতে বদলে যাওয়া}

\subsubsection*{আয়াত ১ — সূরা আল-ফুরকান (২৫) — আয়াত ৭০-৭১}

\subsubsection*{১. সূত্র কার্ড}

\begin{table}[h]\centering\small
\begin{tabular}{ll}
বিষয় & বিবরণ \\
\hline
\textbf{সূরা} & আল-ফুরকান (২৫) — মক্কী \\
\textbf{আয়াত} & ৭০-৭১ \\
\textbf{মূল বিষয়} & তাওবার পর আল্লাহ পাপগুলো নেকিতে বদলে দেন — অতীতের রেকর্ড পরিবর্তন \\
\textbf{প্রসঙ্গ} & শিরক, হত্যা, যিনা — বড় পাপকারীদের তাওবার ঘোষণা \\
\hline
\end{tabular}
\end{table}

\subsubsection*{২. মূল আরবি}

\begin{center}\arab{إِلَّا مَنْ تَابَ وَآمَنَ وَعَمِلَ عَمَلًا صَالِحًا فَأُولَٰئِكَ يُبَدِّلُ اللَّهُ سَيِّئَاتِهِمْ حَسَنَاتٍ ۗ وَكَانَ اللَّهُ غَفُورًا رَحِيمًا}\end{center}

\begin{center}\arab{وَمَنْ تَابَ وَعَمِلَ صَالِحًا فَإِنَّهُ يَتُوبُ إِلَى اللَّهِ مَتَابًا}\end{center}

\subsubsection*{৩. বাংলা অর্থ (\lat{pure})}

\textbf{উচ্চারণ:} ইল্লা- মান তা-বা ওয়া আ-মানা ওয়া 'আমিলা 'আমালান সা-লিহান ফাউলা-ইকা ইউবাদ্দিলুল্লা-হু সাইয়ি-আতিহিম হাসানা-তিন। ওয়া কা-নাল্লা-হু গাফূরার রাহীমা। ওয়া মান তা-বা ওয়া 'আমিলা সা-লিহান ফাইন্নাহু ইয়াতূবু ইলাল্লা-হি মাতা-বা।

\lat{“}কিন্তু যে তাওবা করে, ঈমান আনে এবং সৎকর্ম করে — তাদের পাপগুলো আল্লাহ নেকিতে বদলে দেবেন। আর আল্লাহ ক্ষমাশীল, পরম দয়ালু। আর যে তাওবা করে ও সৎকর্ম করে, সে তো আল্লাহর দিকে সঠিক পথেই ফিরে আসে।\lat{”} (২৫:৭০-৭১)

\subsubsection*{৪. শব্দের ব্যাখ্যা}

\begin{table}[h]\centering\small
\begin{tabular}{ll}
শব্দ & অর্থ \\
\hline
\textbf{তাবা (\arab{تاب})} & তাওবা করা — পাপ থেকে ফেরা \\
\textbf{ইউবাদ্দিলু (\arab{يبدّل})} & বদলে দেন/পরিবর্তন করে দেন — আমলনামায় পাপের জায়গায় নেকি লেখা \\
\textbf{সাইয়ি-আত (\arab{سيئات})} & খারাপ কাজ/পাপ \\
\textbf{মাতাবা (\arab{متابا})} & সঠিক প্রত্যাবর্তন — ফেরার উত্তম পথ \\
\hline
\end{tabular}
\end{table}

\subsubsection*{৫. সংক্ষিপ্ত তাফসীর}

\begin{itemize}
\item \textbf{ইবনে কাসীর:} তাওবা ও সংশোধনের পর আল্লাহ পাপগুলো মুছে দেন এবং সেগুলোর বদলে নেকি লেখা হয়। তাফসীরকারকদের মধ্যে মতভেদ আছে — কেউ বলেন আক্ষরিক অর্থেই আমলনামায় পাপের জায়গায় নেকি লেখা হয় (ইবনে আব্বাস (রাঃ)-এর বিখ্যাত ব্যাখ্যা), কেউ বলেন অর্থ হলো শাস্তির বদলে ক্ষমা ও নেকির প্রতিদান। উভয় মতেই মূল বিষয় এক — \textbf{তাওবার পর পাপের শাস্তি ও হিসাবের ভয় থাকে না; বরং ক্ষমা ও কল্যাণ।}
\item \textbf{ইবনে রজব (রহ.)} বলেন: যে ব্যক্তি পাপের জন্য সত্যিই অনুতপ্ত হয়, তার প্রতিটি পাপ তাকে নেক কাজের দিকে ঠেলে দেয় — ফলে পাপগুলোই নেকিতে রূপান্তরিত হয়। ইবনে কাসীর আরও বলেন — এই আয়াতের উপর ভিত্তি করেই আলেমগণ বলেন, \textbf{তাওবাকারীকে তার পাপের জন্য কিয়ামতে জিজ্ঞাসাবাদ করা হবে না, কারণ তাওবা সেগুলো মুছে দিয়েছে।}
\end{itemize}
\subsubsection*{৬. সংশ্লিষ্ট দলিল}

\begin{itemize}
\item \textbf{হাদিস (বুখারি ৭৫০৭; মুসলিম ২৭৫৮):} বান্দা বারবার পাপ করে \lat{“}রব, আমাকে ক্ষমা করো\lat{”} বলে — আল্লাহ প্রতিবার বলেন: \lat{“}আমার বান্দা জানে তার এমন রব আছে যিনি পাপ ক্ষমা করেন এবং পাপের শাস্তিও দিতে পারেন — আমি আমার বান্দাকে ক্ষমা করে দিলাম।\lat{”}
\item \textbf{হাদিস (মুসলিম ২৭৬৮):} নাজওয়া হাদিস — ক্ষমার পর তাকে \lat{“}নেকির আমলনামা\lat{”} দেওয়া হয়।
\end{itemize}
\subsubsection*{আয়াত ২ — সূরা হুদ (১১) — আয়াত ১১৪}

\subsubsection*{১. সূত্র কার্ড}

\begin{table}[h]\centering\small
\begin{tabular}{ll}
বিষয় & বিবরণ \\
\hline
\textbf{সূরা} & হুদ (১১) — মক্কী \\
\textbf{আয়াত} & ১১৪ \\
\textbf{মূল বিষয়} & ভালো কাজ মন্দ কাজ মুছে দেয় — পাপের প্রভাব দূর করার পথ \\
\textbf{প্রসঙ্গ} & নামাজের গুরুত্বের প্রসঙ্গে \\
\hline
\end{tabular}
\end{table}

\subsubsection*{২. মূল আরবি}

\begin{center}\arab{وَأَقِمِ الصَّلَاةَ طَرَفَيِ النَّهَارِ وَزُلَفًا مِنَ اللَّيْلِ ۚ إِنَّ الْحَسَنَاتِ يُذْهِبْنَ السَّيِّئَاتِ ۚ ذَٰلِكَ ذِكْرَىٰ لِلذَّاكِرِينَ}\end{center}

\subsubsection*{৩. বাংলা অর্থ (\lat{pure})}

\textbf{উচ্চারণ:} ওয়া আক্বিমিস সালাতা ত্বারাফাইন নাহারি ওয়া যুলাফাম মিনাল লাইল। ইন্নাল হাসানা-তি ইউযহিবনাস সাইয়ি-আ-তি। যা-লিকা যিকরা- লিয যা-কিরীন।

\lat{“}আর দিনের দুই প্রান্তে ও রাতের কিছু অংশে নামাজ কায়েম করো। নিশ্চয়ই ভালো কাজ মন্দ কাজ মুছে দেয়। এটি স্মরণকারীদের জন্য উপদেশ।\lat{”} (১১:১১৪)

\subsubsection*{৪. শব্দের ব্যাখ্যা}

\begin{table}[h]\centering\small
\begin{tabular}{ll}
শব্দ & অর্থ \\
\hline
\textbf{ইউযহিবনা (\arab{يذهبن})} & মুছে দেয়/দূর করে দেয় \\
\textbf{আল-হাসানাত} & ভালো কাজ — বিশেষত নামাজ (প্রসঙ্গ অনুযায়ী) \\
\textbf{আস-সাইয়ি-আত} & মন্দ কাজ/পাপ \\
\hline
\end{tabular}
\end{table}

\subsubsection*{৫. সংক্ষিপ্ত তাফসীর}

\begin{itemize}
\item \textbf{ইবনে কাসীর:} এই আয়াত নাজিলের কারণ — ইবনে মাসউদ (রাঃ) থেকে বর্ণিত: এক ব্যক্তি একজন নারীকে চুমু দিয়েছিল, নবীজির (সা.) কাছে এসে তা জানাল; তখন এই আয়াত নাজিল হলো। লোকটি জিজ্ঞেস করল: \lat{“}এটি কি শুধু আমার জন্য?\lat{”} নবীজি (সা.) বললেন: \lat{“}বরং আমার উম্মতের মধ্যে যে এ অনুযায়ী আমল করবে — তার জন্যই এটি।\lat{”} (বুখারি ৪৬৮৭; মুসলিম ২৭৬৩)
\item অর্থাৎ \textbf{পাপ হয়ে গেলে তার পরের ভালো কাজগুলো পাপের প্রভাব দূর করে দেয়} — পাপ যেন আমলনামায় থাকেই না। এটি কিয়ামতের হিসাবকে সহজ করে।
\end{itemize}
\subsubsection*{৬. সংশ্লিষ্ট দলিল}

\begin{itemize}
\item \textbf{হাদিস (বুখারি ৪৬৮৭; মুসলিম ২৭৬৩):} চুমুর পাপের ঘটনা — \lat{“}বরং আমার উম্মতের মধ্যে যে এ অনুযায়ী আমল করবে, তার জন্যই এটি।\lat{”}
\item \textbf{আয়াত (৪:৩১):} \lat{“}তোমরা বড় পাপগুলো থেকে বিরত থাকলে আমি তোমাদের ছোটখাটো পাপ মোচন করব।\lat{”}
\end{itemize}
\subsubsection*{আয়াত ৩ — সূরা আন-নিসা (৪) — আয়াত ৩১}

\subsubsection*{১. সূত্র কার্ড}

\begin{table}[h]\centering\small
\begin{tabular}{ll}
বিষয় & বিবরণ \\
\hline
\textbf{সূরা} & আন-নিসা (৪) — মাদানী \\
\textbf{আয়াত} & ৩১ \\
\textbf{মূল বিষয়} & বড় পাপ পরিহার করলে ছোট পাপ মোচন — ক্ষমার প্রশস্ত দরজা \\
\textbf{প্রসঙ্গ} & হারাম উপার্জন ও অন্যায়-অধিকার নিষেধের পর \\
\hline
\end{tabular}
\end{table}

\subsubsection*{২. মূল আরবি}

\begin{center}\arab{إِنْ تَجْتَنِبُوا كَبَائِرَ مَا تُنْهَوْنَ عَنْهُ نُكَفِّرْ عَنْكُمْ سَيِّئَاتِكُمْ وَنُدْخِلْكُمْ مُدْخَلًا كَرِيمًا}\end{center}

\subsubsection*{৩. বাংলা অর্থ (\lat{pure})}

\textbf{উচ্চারণ:} ইন তাজতানিবূ কাবা-ইরা মা- তুনহাওনা 'আনহু নুকাফফির 'আনকুম সাইয়ি-আতিকুম ওয়া নুদখিলকুম মুদখালান কারীমা।

\lat{“}তোমরা যে বড় পাপগুলো থেকে নিষেধ হয়েছ, সেগুলো থেকে বিরত থাকলে আমি তোমাদের ছোটখাটো পাপ মোচন করে দেব এবং তোমাদেরকে সম্মানজনক স্থানে প্রবেশ করাব।\lat{”} (৪:৩১)

\subsubsection*{৪. শব্দের ব্যাখ্যা}

\begin{table}[h]\centering\small
\begin{tabular}{ll}
শব্দ & অর্থ \\
\hline
\textbf{কাবা-ইর (\arab{كبائر})} & বড় পাপসমূহ \\
\textbf{নুকাফফির (\arab{نكفّر})} & মোচন করি/মুছে দিই — ঢেকে দিই \\
\textbf{মুদখালান কারিমা} & সম্মানজনক প্রবেশস্থল — জান্নাত \\
\hline
\end{tabular}
\end{table}

\subsubsection*{৫. সংক্ষিপ্ত তাফসীর}

\begin{itemize}
\item \textbf{ইবনে কাসীর:} এই আয়াত তাওবাকারীদের জন্য আশার বার্তা — যে ব্যক্তি বড় পাপ থেকে বেঁচে থাকে, আল্লাহ তার ছোট পাপ মোচন করেন। এর অর্থ এই নয় যে বড় পাপ করলে আল্লাহ ক্ষমা করেন না — বরং তাওবার শর্তে বড় পাপও ক্ষমা হয় (২৫:৭০, ৩৯:৫৩)। এখানে সাধারণ নিয়ম বলা হচ্ছে।
\end{itemize}
\subsubsection*{৬. সংশ্লিষ্ট দলিল}

\begin{itemize}
\item \textbf{আয়াত (আল-ফুরকান ২৫:৭০):} তাওবার পর বড় পাপও নেকিতে বদলে যায়।
\item \textbf{হাদিস (বুখারি ৬৫৩৬; মুসলিম ২৮৭৬):} \lat{“}যার হিসাব পুঙ্খানুপুঙ্খভাবে নেওয়া হবে, সে শাস্তি পাবে\lat{”} — অর্থাৎ সহজ হিসাব মানে হিসাবই নেই, বরং আমলের পেশকরণ (আরয)।
\end{itemize}
\medskip\hrule\medskip

\subsection*{থিম ২ — কিয়ামতের দিন আমলনামা ও হিসাবের বাস্তবতা}

\subsubsection*{আয়াত ৪ — সূরা আল-কাহফ (১৮) — আয়াত ৪৯}

\subsubsection*{১. সূত্র কার্ড}

\begin{table}[h]\centering\small
\begin{tabular}{ll}
বিষয় & বিবরণ \\
\hline
\textbf{সূরা} & আল-কাহফ (১৮) — মক্কী \\
\textbf{আয়াত} & ৪৯ \\
\textbf{মূল বিষয়} & আমলনামা পেশ — ছোট-বড় সব লেখা থাকে; কিন্তু তাওবার পর ক্ষমা \\
\textbf{প্রসঙ্গ} & কিয়ামতের ভয়াবহতার বর্ণনা — মুজরিমদের প্রতিক্রিয়া \\
\hline
\end{tabular}
\end{table}

\subsubsection*{২. মূল আরবি}

\begin{center}\arab{وَوُضِعَ الْكِتَابُ فَتَرَى الْمُجْرِمِينَ مُشْفِقِينَ مِمَّا فِيهِ وَيَقُولُونَ يَا وَيْلَتَنَا مَالِ هَٰذَا الْكِتَابِ لَا يُغَادِرُ صَغِيرَةً وَلَا كَبِيرَةً إِلَّا أَحْصَاهَا ۚ وَوَجَدُوا مَا عَمِلُوا حَاضِرًا ۗ وَلَا يَظْلِمُ رَبُّكَ أَحَدًا}\end{center}

\subsubsection*{৩. বাংলা অর্থ (\lat{pure})}

\textbf{উচ্চারণ:} ওয়া ওয়ুদি'আল কিতা-বু ফাতারাল মুজরিমীনা মুশফিক্বীনা মিম্মা- ফীহি ওয়া ইয়াক্বুলূনা ইয়া- ওয়াইলাতানা- মা- লি হা-যাল কিতা-বি লা- ইউগা-দিরু সাগীরাতান ওয়ালা- কাবীরাতান ইল্লা- আহসা-হা-। ওয়া ওয়াজাদূ মা- 'আমিলূ হা-দিরা-। ওয়ালা- ইয়াজলিমু রব্বুকা আহাদা।

\lat{“}আর আমলনামা রাখা হবে। তুমি দেখবে অপরাধীরা তাতে যা আছে তা নিয়ে ভীত-সন্ত্রস্ত; তারা বলবে: 'হায় আমাদের দুর্ভাগ্য! এ কেমন আমলনামা! এ তো ছোট-বড় কোনো কিছুই বাদ দেয়নি, সবই লিপিবদ্ধ করেছে!' তারা যা করেছে সবই পাবে — তোমার রব কারো প্রতি জুলুম করেন না।\lat{”} (১৮:৪৯)

\subsubsection*{৪. শব্দের ব্যাখ্যা}

\begin{table}[h]\centering\small
\begin{tabular}{ll}
শব্দ & অর্থ \\
\hline
\textbf{আল-মুজরিমীন} & অপরাধীরা — যারা তাওবা করেনি \\
\textbf{লা ইউগাদিরু} & বাদ দেয় না — সব লিপিবদ্ধ \\
\textbf{আহসাহা} & গণনা করেছে/লিপিবদ্ধ করেছে \\
\textbf{লা ইয়াজলিমু} & জুলুম করেন না — পূর্ণ ন্যায়বিচার \\
\hline
\end{tabular}
\end{table}

\subsubsection*{৫. সংক্ষিপ্ত তাফসীর}

\begin{itemize}
\item \textbf{ইবনে কাসীর:} এই আয়াত বলে — আমলনামায় সব কাজ লেখা থাকে, কিছু বাদ যায় না। তবে এই ভয়-ভীতি ও আফসোস \textbf{মুজরিম (অপরাধী)দের} — যারা তাওবা করেনি। তাওবাকারীর ক্ষেত্রে কুরআনই বলে (২৫:৭০) পাপগুলো নেকিতে বদলে দেওয়া হয়, এবং নাজওয়া হাদিসে (বুখারি ২৪৪১) ক্ষমার পর তাকে \textbf{নেকির আমলনামা} দেওয়া হয়।
\item অর্থাৎ: \textbf{যে পাপ আল্লাহ ক্ষমা করেছেন, তা আর হিসাবের বোঝা নয়} — আমলনামার যে অংশে ছিল, সেটি ক্ষমা/নেকি দ্বারা প্রতিস্থাপিত। ১৮:৪৯ সেই ব্যক্তির কথা, যে তাওবা ছাড়াই আল্লাহর সামনে হাজির হয়।
\item \textbf{সা'দী:} আমলনামার পূর্ণতা আল্লাহর ন্যায়বিচারের প্রমাণ — কারো কোনো কাজের হিসাব হারিয়ে যাবে না; কিন্তু এটিই আবার ক্ষমার ক্ষেত্রে আল্লাহর অনুগ্রহের মুখ নয় — অনুগ্রহের দরজা তাওবার মাধ্যমে খোলা।
\end{itemize}
\subsubsection*{৬. সংশ্লিষ্ট দলিল}

\begin{itemize}
\item \textbf{হাদিস (বুখারি ২৪৪১; মুসলিম ২৭৬৮):} নাজওয়া — মুমিনের পাপ ক্ষমা করে তাকে নেকির আমলনামা দেওয়া হয়।
\item \textbf{আয়াত (আল-ফুরকান ২৫:৭০):} তাওবাকারীর পাপ নেকিতে বদল।
\end{itemize}
\subsubsection*{আয়াত ৫ — সূরা আল-হাক্কাহ (৬৯) — আয়াত ১৮-২৫}

\subsubsection*{১. সূত্র কার্ড}

\begin{table}[h]\centering\small
\begin{tabular}{ll}
বিষয় & বিবরণ \\
\hline
\textbf{সূরা} & আল-হাক্কাহ (৬৯) — মক্কী \\
\textbf{আয়াত} & ১৮-২৫ \\
\textbf{মূল বিষয়} & কিয়ামতে সবার আমলনামা পেশ — নেককারের আনন্দ, বদকারের আফসোস \\
\textbf{প্রসঙ্গ} & কিয়ামতের দৃশ্য — হিসাবের প্রকৃতি \\
\hline
\end{tabular}
\end{table}

\subsubsection*{২. মূল আরবি}

\begin{center}\arab{يَوْمَئِذٍ تُعْرَضُونَ لَا تَخْفَىٰ مِنْكُمْ خَافِيَةٌ}\end{center}

\begin{center}\arab{فَأَمَّا مَنْ أُوتِيَ كِتَابَهُ بِيَمِينِهِ فَيَقُولُ هَاؤُمُ اقْرَءُوا كِتَابِيَهْ}\end{center}

\subsubsection*{৩. বাংলা অর্থ (\lat{pure})}

\textbf{উচ্চারণ:} ইয়াওমাইযিন তু'রাদূনা লা- তাখফা- মিনকুম খা-ফিয়াহ। ফাআম্মা- মান উতিয়া কিতাবাহু বিয়ামীনিহী ফাইয়াক্বুলু হা-উমুক্বরাউ কিতাবিয়াহ।

\lat{“}সেদিন তোমাদেরকে পেশ করা হবে — তোমাদের কোনো গোপন বিষয়ই লুকানো থাকবে না। তখন যাকে তার আমলনামা ডান হাতে দেওয়া হবে, সে বলবে: 'নাও, আমার আমলনামা পড়ো!'\lat{”} (৬৯:১৮-১৯)

\subsubsection*{৪. শব্দের ব্যাখ্যা}

\begin{table}[h]\centering\small
\begin{tabular}{ll}
শব্দ & অর্থ \\
\hline
\textbf{তু'রাদূন} & পেশ করা হবে — হিসাবের জন্য উপস্থাপন \\
\textbf{লা তাখফা খাফিয়াহ} & কোনো গোপন বিষয় লুকানো থাকবে না — সব প্রকাশ \\
\textbf{বিয়ামীনিহি} & তার ডান হাতে — সৌভাগ্যের চিহ্ন \\
\hline
\end{tabular}
\end{table}

\subsubsection*{৫. সংক্ষিপ্ত তাফসীর}

\begin{itemize}
\item \textbf{ইবনে কাসীর:} ৬৯:১৮-তে \lat{“}কোনো গোপন বিষয় লুকানো থাকবে না\lat{”} — অর্থ আল্লাহর কাছে সব প্রকাশ; হিসাব নেওয়া হবে। তবে মুমিনের জন্য এই পেশকরণ \textbf{আরয (প্রদর্শন)} — মুসলিম ২৮৭৬-এর ব্যাখ্যা অনুযায়ী \lat{“}সহজ হিসাব\lat{”} মানে পুঙ্খানুপুঙ্খ জিজ্ঞাসাবাদ নয়, বরং আমলের পেশকরণ।
\item \textbf{ইবনে কাসীর:} ডান হাতে আমলনামা পাওয়া = জান্নাতের সুসংবাদ। ক্ষমা করা পাপ এখানে আর বোঝা নয় — কারণ তাওবা সেগুলো মুছে দিয়েছে (২৫:৭০)।
\end{itemize}
\subsubsection*{৬. সংশ্লিষ্ট দলিল}

\begin{itemize}
\item \textbf{হাদিস (বুখারি ৬৫৩৬; মুসলিম ২৮৭৬):} \lat{“}সহজ হিসাব\lat{”} (৮৪:৮) মানে আরয — প্রদর্শন; পুঙ্খানুপুঙ্খ হিসাব মানেই শাস্তি।
\item \textbf{আয়াত (আল-ইনশিকাক ৮৪:৭-৮):} ডান হাতে আমলনামা পেলে সহজ হিসাব।
\end{itemize}
\subsubsection*{আয়াত ৬ — সূরা আত-তারিক (৮৬) — আয়াত ৯}

\subsubsection*{১. সূত্র কার্ড}

\begin{table}[h]\centering\small
\begin{tabular}{ll}
বিষয় & বিবরণ \\
\hline
\textbf{সূরা} & আত-তারিক (৮৬) — মক্কী \\
\textbf{আয়াত} & ৯ \\
\textbf{মূল বিষয়} & যেদিন গোপন বিষয় প্রকাশিত হবে — কিন্তু মুমিনের ক্ষেত্রে গোপন রাখার প্রতিশ্রুতি \\
\textbf{প্রসঙ্গ} & কিয়ামতের শপথ — মানুষের প্রত্যাবর্তন \\
\hline
\end{tabular}
\end{table}

\subsubsection*{২. মূল আরবি}

\begin{center}\arab{يَوْمَ تُبْلَى السَّرَائِرُ}\end{center}

\subsubsection*{৩. বাংলা অর্থ (\lat{pure})}

\textbf{উচ্চারণ:} ইয়াওমা তুবলাস সারা-ইরু।

\lat{“}সেদিন গোপন বিষয়সমূহ পরীক্ষা করা হবে।\lat{”} (৮৬:৯)

\subsubsection*{৪. শব্দের ব্যাখ্যা}

\begin{table}[h]\centering\small
\begin{tabular}{ll}
শব্দ & অর্থ \\
\hline
\textbf{তুবলা (\arab{تبلى})} & পরীক্ষা করা হবে/প্রকাশ করা হবে \\
\textbf{আস-সারাইর} & গোপন বিষয়সমূহ — অন্তরের গোপন \\
\hline
\end{tabular}
\end{table}

\subsubsection*{৫. সংক্ষিপ্ত তাফসীর}

\begin{itemize}
\item \textbf{ইবনে কাসীর:} গোপন বিষয় প্রকাশ — অর্থ আল্লাহর কাছে সব প্রকাশ হয়ে যাবে; হিসাবের দিন লুকানোর কিছু থাকবে না।
\item \textbf{গুরুত্বপূর্ণ পয়েন্ট:} এই প্রকাশ মানে এই নয় যে সবাই সবার পাপ দেখবে। নাজওয়া হাদিসে (বুখারি ২৪৪১) স্পষ্ট — মুমিনকে আল্লাহ কাফি (আড়াল) দিয়ে একান্তে তাঁর পাপ স্মরণ করিয়ে ক্ষমা করেন; \textbf{সবার সামনে প্রকাশ হয় কাফির-মুনাফিকদের।} ইবনে হাজার (রহ.) বলেন — হাদিসটি মুমিনদের পাপ কিয়ামতে গোপন রাখার দলিল; কারণ কাফির-মুনাফিকদের ক্ষেত্রে এই গোপন রাখার কথা বাদ দেওয়া হয়েছে।
\item \textbf{সা'দী:} প্রকাশের ভয় তাওবাকারীকে স্পর্শ করে না — যার পাপ ক্ষমা হয়ে গেছে, তার গোপন বিষয় আল্লাহ নিজেই ঢেকে রাখেন।
\end{itemize}
\subsubsection*{৬. সংশ্লিষ্ট দলিল}

\begin{itemize}
\item \textbf{হাদিস (বুখারি ২৪৪১; মুসলিম ২৭৬৮):} মুমিনের পাপ একান্তে স্মরণ করিয়ে ক্ষমা — কাফির-মুনাফিকের প্রকাশ।
\item \textbf{হাদিস (বুখারি ৬০৬৯):} \lat{“}আমার সব উম্মত ক্ষমা পাবে, কিন্তু মুজাহির (যে নিজের পাপ প্রকাশ করে) ছাড়া।\lat{”}
\end{itemize}
\medskip\hrule\medskip

\subsection*{থিম ৩ — ক্ষমার সীমা: শিরক, মানুষের হক ও গোপন করার বিধান}

\subsubsection*{আয়াত ৭ — সূরা আন-নিসা (৪) — আয়াত ৪৮ ও ১১৬}

\subsubsection*{১. সূত্র কার্ড}

\begin{table}[h]\centering\small
\begin{tabular}{ll}
বিষয় & বিবরণ \\
\hline
\textbf{সূরা} & আন-নিসা (৪) — মাদানী \\
\textbf{আয়াত} & ৪৮ ও ১১৬ \\
\textbf{মূল বিষয়} & শিরক ছাড়া সব পাপ — আল্লাহ চাইলে ক্ষমা করেন; তাওবার পর শিরকও ক্ষমা হয় \\
\textbf{প্রসঙ্গ} & ক্ষমার প্রশস্ততার সাথে সতর্কতার ভারসাম্য \\
\hline
\end{tabular}
\end{table}

\subsubsection*{২. মূল আরবি}

\begin{center}\arab{إِنَّ اللَّهَ لَا يَغْفِرُ أَنْ يُشْرَكَ بِهِ وَيَغْفِرُ مَا دُونَ ذَٰلِكَ لِمَنْ يَشَاءُ ۚ وَمَنْ يُشْرِكْ بِاللَّهِ فَقَدِ افْتَرَىٰ إِثْمًا عَظِيمًا}\end{center}

\subsubsection*{৩. বাংলা অর্থ (\lat{pure})}

\textbf{উচ্চারণ:} ইন্নাল্লা-হা লা- ইয়াগফিরু আই ইউশরাকা বিহী ওয়া ইয়াগফিরু মা- দূনা যা-লিকা লিমাই ইয়াশা-উ। ওয়া মাই ইউশরিক বিল্লা-হি ফাক্বাদিফতার ইসমান 'আযীমা।

\lat{“}নিশ্চয়ই আল্লাহ শিরক ক্ষমা করেন না; কিন্তু এছাড়া যাকে ইচ্ছা ক্ষমা করেন। আর যে আল্লাহর সাথে শিরক করে, সে তো মহাপাপ রচনা করল।\lat{”} (৪:৪৮)

\subsubsection*{৪. শব্দের ব্যাখ্যা}

\begin{table}[h]\centering\small
\begin{tabular}{ll}
শব্দ & অর্থ \\
\hline
\textbf{লা ইয়াগফিরু} & ক্ষমা করেন না — শিরকে মৃত্যু পর্যন্ত অটল থাকলে \\
\textbf{মা দূনা যালিকা} & এছাড়া যা কিছু — শিরক ছাড়া সব \\
\textbf{লিমাই ইয়াশা} & যাকে ইচ্ছা — আল্লাহর ইচ্ছাধীন \\
\hline
\end{tabular}
\end{table}

\subsubsection*{৫. সংক্ষিপ্ত তাফসীর}

\begin{itemize}
\item \textbf{ইবনে কাসীর:} এই আয়াত বলে — শিরকে মৃত্যুবরণকারী ছাড়া বাকি সব পাপ আল্লাহর ইচ্ছাধীন; তিনি চাইলে ক্ষমা করেন। কিন্তু \textbf{তাওবার পর শিরকও ক্ষমা হয়} — ৩৯:৫৩, ২৫:৭০-এ তার ঘোষণা (ওয়াশশি ও মুশরিকদের তাওবার ঘটনা — ইবনে কাসীর)।
\item \textbf{মূল নীতি:} ক্ষমার প্রশস্ততা অসীম — তবে শিরকে মৃত্যু ও মানুষের হক নষ্ট (অন্যায়ভাবে) ক্ষমার সাধারণ বিধানের বাইরে; দ্বিতীয়টি ফেরত দেওয়া/ক্ষমা আদায় করা ছাড়া।
\item \textbf{সা'দী:} ক্ষমা আল্লাহর হাতে — তাওবাকারী নিশ্চিন্ত হতে পারে, কিন্তু শিরক ও জুলুমের জন্য সাবধান থাকা জরুরি।
\end{itemize}
\subsubsection*{৬. সংশ্লিষ্ট দলিল}

\begin{itemize}
\item \textbf{আয়াত (আয-যুমার ৩৯:৫৩):} \lat{“}তোমরা আল্লাহর রহমত থেকে নিরাশ হয়ো না — আল্লাহ সব গুনাহ ক্ষমা করেন।\lat{”}
\item \textbf{হাদিস (বুখারি ২৪৪১; মুসলিম ২৭৬৮):} নাজওয়া — ক্ষমার পর নেকির আমলনামা।
\end{itemize}
\medskip\hrule\medskip

\subsection*{থিম ৪ — হিসাবের দুই রূপ ও সহজ হিসাবের অর্থ}

\subsubsection*{আয়াত ৮ — সূরা আল-ইনশিকাক (৮৪) — আয়াত ৭-৮}

\subsubsection*{১. সূত্র কার্ড}

\begin{table}[h]\centering\small
\begin{tabular}{ll}
বিষয় & বিবরণ \\
\hline
\textbf{সূরা} & আল-ইনশিকাক (৮৪) — মক্কী \\
\textbf{আয়াত} & ৭-৮ \\
\textbf{মূল বিষয়} & ডান হাতে আমলনামা — সহজ হিসাব \\
\textbf{প্রসঙ্গ} & কিয়ামতের দৃশ্য — আমলনামার বিতরণ \\
\hline
\end{tabular}
\end{table}

\subsubsection*{২. মূল আরবি}

\begin{center}\arab{فَأَمَّا مَنْ أُوتِيَ كِتَابَهُ بِيَمِينِهِ ۝ فَسَوْفَ يُحَاسَبُ حِسَابًا يَسِيرًا}\end{center}

\subsubsection*{৩. বাংলা অর্থ (\lat{pure})}

\textbf{উচ্চারণ:} ফাআম্মা- মান উতিয়া কিতাবাহু বিয়ামীনিহী, ফাসাওফা ইউহাসাবু হিসাবান ইয়াসীরা।

\lat{“}অতঃপর যাকে তার আমলনামা ডান হাতে দেওয়া হবে, সে তো সহজ হিসাবে হিসাব পেশ করবে।\lat{”} (৮৪:৭-৮)

\subsubsection*{৪. শব্দের ব্যাখ্যা}

\begin{table}[h]\centering\small
\begin{tabular}{ll}
শব্দ & অর্থ \\
\hline
\textbf{উতিয়া বিয়ামীনিহি} & ডান হাতে দেওয়া হবে — সৌভাগ্যবান \\
\textbf{হিসাবান ইয়াসিরা} & সহজ হিসাব — পুঙ্খানুপুঙ্খ জিজ্ঞাসাবাদ ছাড়া \\
\hline
\end{tabular}
\end{table}

\subsubsection*{৫. সংক্ষিপ্ত তাফসীর}

\begin{itemize}
\item \textbf{ইবনে কাসীর:} সহজ হিসাবের ব্যাখ্যা নবীজি (সা.) নিজেই করেছেন — আয়েশা (রা.) জিজ্ঞেস করলেন, \lat{“}সহজ হিসাব\lat{”} কী? নবীজি বললেন: \lat{“}সেটি তো আরয (আমলের পেশকরণ); আর যার হিসাব পুঙ্খানুপুঙ্খভাবে নেওয়া হবে, সে ধ্বংস হয়ে যাবে।\lat{”} (বুখারি ৬৫৩৬; মুসলিম ২৮৭৬)
\item অর্থাৎ \textbf{তাওবাকারীর পাপগুলো আর পুঙ্খানুপুঙ্খ জিজ্ঞাসাবাদের বিষয় নয়} — তার হিসাব সহজ, কারণ পাপ ক্ষমা হয়ে গেছে (২৫:৭০)।
\end{itemize}
\subsubsection*{৬. সংশ্লিষ্ট দলিল}

\begin{itemize}
\item \textbf{হাদিস (বুখারি ৬৫৩৬; মুসলিম ২৮৭৬):} আয়েশা (রা.)-এর প্রশ্ন ও নবীজির উত্তর — \lat{“}সেটি আরয; পুঙ্খানুপুঙ্খ হিসাব মানেই শাস্তি।\lat{”}
\item \textbf{আয়াত (আল-হাক্কাহ ৬৯:১৮-২৫):} ডান হাতে আমলনামা — নেককারের আনন্দ।
\end{itemize}
\medskip\hrule\medskip

\subsection*{মূল সিদ্ধান্ত (দলিলের আলোকে)}

\begin{enumerate}
\item \textbf{তাওবার পর পাপ মুছে যায় / নেকিতে বদলে যায়} — ২৫:৭০; ১১:১১৪; ৪:৩১।
\item \textbf{ক্ষমা করা পাপ আর পুঙ্খানুপুঙ্খ হিসাবে আসে না} — ৮৪:৭-৮-এর ব্যাখ্যা (বুখারি ৬৫৩৬): সহজ হিসাব = আরয।
\item \textbf{মুমিনের পাপ কিয়ামতে প্রকাশ হয় না — একান্তে ক্ষমা} — নাজওয়া হাদিস (বুখারি ২৪৪১; মুসলিম ২৭৬৮); ৮৬:৯-এর প্রকৃত অর্থ আল্লাহর কাছে সব প্রকাশ, সবার সামনে নয়।
\item \textbf{শিরকে মৃত্যু ক্ষমার বাইরে; বাকি সব তাওবার শর্তে ক্ষমা} — ৪:৪৮; ৩৯:৫৩।
\item \textbf{মানুষের হক নষ্টের হিসাব ভিন্ন} — ফেরত/ক্ষমা আদায় ছাড়া ক্ষমা প্রযোজ্য নয় (মুসলিম ২৫৮১ — মুফলিস হাদিস)।
\end{enumerate}

\end{document}
